\documentclass[letterpaper]{elsarticle}
\usepackage[utf8]{inputenc}
\usepackage[spanish]{babel}
\usepackage{amsmath}
\usepackage{graphicx}
\usepackage{comment}
\usepackage{tikz}
\usepackage{subfig}
\usepackage{SIunits}
\usepackage{pgfgantt}%Gantt
\usepackage{transparent}%marca de agua
\usepackage{tcolorbox}%frames:
\usepackage{eso-pic}%marca de agua
\usepackage{pgfplots}
\pgfplotsset{width=2.5in,compat=1.10, every axis legend/.append style={
        at={(0.5,-0.52)},
        anchor=south}}
\usepackage[colorinlistoftodos]{todonotes}
\newcommand{\oh}{\emph{OH}} %defines an \oh command 
%----------------------
\newcommand{\imref}[1]{\figurename \ref{#1}}
%-------------------
\renewcommand\thesection{\arabic{section}}
\renewcommand\thesubsection{\thesection.\arabic{subsection}}
%-----------------
%Encabezado y pie de página:
\usepackage{fancyhdr}
\pagestyle{fancy}

\lhead{\begin{picture}(0,0)
\put(0,10){
\includegraphics[width=4cm]{SEP}}		
\end{picture}}


\chead{\tiny \textsc{Informe técnico de avance para proyectos financiados por el Tecnológico Nacional de México}}
\rhead{
\tiny TECNOLÓGICO NACIONAL DE MÉXICO\\
Coordinación Sectorial Académica\\
Dirección de Estudios de Posgrado e Investigación\\
\bigskip}
\lfoot{CI-PY-ITA-01/ITM}
\renewcommand{\headrulewidth}{0pt}
\headsep=35pt
%------------------------
%Marca de agua
\newcommand\BackgroundPic{
\put(0,0){\parbox[b][\paperheight]{\paperwidth}{%
\vfill
\centering
{\includegraphics[width=6in,
keepaspectratio]{escudoNacionalt}}%
\vfill
}}}
%-----------------

%Datos del documento:
\def\myTitle{Título del proyecto}



\def\myAuthor{Chávez-Campos, Gerardo~Marx}

\def\myInstitution{Instituto Tecnológico de Morelia}
\def\myProject{2234-5}

\newcommand{\declaracion}{Antes de proporcionar la información solicitada, lea cuidadosamente cada uno de los rubros que contiene el presente formato. Es importarte no modificarlo.}



%.------------------------------
%Body of document:
%------------------------:)
\begin{document}
\AddToShipoutPicture*{\BackgroundPic}
\declaracion
\begin{flushright}
	\textbf\today
\end{flushright}
\section{Identificación del proyecto}
\begin{center}
\begin{tcolorbox}[title=Institución:, colback=white]
		\centering\textbf\myInstitution	
\end{tcolorbox}
\bigskip
\tcbox[title=Director del proyecto:, colback=white]{\centering\textbf\myAuthor}
\end{center}

\noindent \textbf{Clave del proyecto:} \myProject\par\bigskip
\noindent \textbf{Título:} \myTitle\bigskip

\noindent \textbf{Nivel académico:} ( ) Licenciatura ( ) Maestría\smallskip

\noindent \textbf{Nombre del posgrado donde se realiza el proyecto:} Maestría en ciencias en ingeniería electrónica\smallskip

\noindent \textbf{Nombre de la licenciatura donde se realiza el proyecto}: \smallskip

\noindent \textbf{Línea de investigación o de trabajo}: LGAC \smallskip

\noindent \textbf{Tipo de investigación:}
\begin{center}
Básica ( ) Aplicada ( ) Desarrollo tecnológico ( )
\end{center} 

\noindent \textbf{Duración del proyecto}: ( ) 1 año ( ) 2 Años \par\smallbreak

\noindent \textbf{Fecha de inicio del proyecto:} 02 de mayo 2020 \textbf{Fecha de termino del proyecto:} 02 de mayo 2027.\smallbreak

\noindent \textbf{Periodo que reporta:} 01 noviembre 2015 al 30 de noviembre 2015\smallskip

\begin{center}
	\textbf{No. report:}  (x) 1er Avance ( ) 2do Avance ( ) 3er Avance \smallbreak
\end{center} 
%----------------------------
\section{Resultados}
\subsection{Resumen del avance del proyecto} 
%Describa brevemente las actividades más relevantes realizadas y los resultados obtenidos en esta etapa, de tal forma que permita evaluar el cumplimiento del programa de trabajo. 


\subsection{Avance acumulado}
%Indique en términos porcentuales el avance acumulado del proyecto, con base en las metas alcanzadas y planteadas en el protocolo.
\subsection{Avance de objetivos}
%Indique el estado de avance en el logro de los objetivos planteados en el proyecto (máximo media cuartilla).

\subsection{Avance de metas}
%

\begin{table}[!h]
\renewcommand{\arraystretch}{1.3}
\caption{\textsc{Lista de metas del periodo}}
\label{tb:metas}
\centering
\begin{tabular}{lccc}
\bfseries  & \bfseries cantidad&\bfseries cantidad&\bfseries cantidad en\\
\bfseries Metas cuantificables & \bfseries programada &\bfseries logradas& \bfseries proceso\\
\hline
Residencias profesionales & 800& 40.03 & 185.07\\
Tesis de licenciatura &2& 0 & 2\\
Trabajo de especialidad& -& - & -\\
Tesis de maestría& -& - & -\\
Tesis de doctorado& -& - & -\\
Publicaciones& - & - & 1\\
Patentes& - & - & 1\\
Prototipos& - & - & 1\\
Asistencia a eventos académicos &- &-&-\\
\hline
\end{tabular}
\end{table}

\subsection{Metodología empleada}
%1.	Describa la metodología empleada en esta etapa para el alcance de los objetivos (máximo dos cuartillas).

\subsection{Actividades realizadas}
%1.	Enliste las actividades realizadas durante este periodo, indicando el porcentaje de avance.
\begin{table}[h!]
\begin{center}
\caption{Actividades realizadas en el periodo reportado}

\begin{tikzpicture}
\begin{ganttchart}[
x unit =1.7cm, 
y unit title=0.6cm, 
y unit chart=1cm,
vgrid,
progress=today,
today=2015-11,
%inline,
time slot format=isodate-yearmonth, 
compress calendar
]
{2015-05}%fecha de inicio
{2015-11}%fecha final
\gantttitlecalendar{year, month=name} \\
\ganttbar{Tarea 1}{2015-05}{2015-06}%actividad 1
\end{ganttchart}
\end{tikzpicture}
\end{center}
\end{table}

\subsection{Participación e integración del grupo de trabajo}

%Describa la participación e integración del grupo de trabajo. Indique en la siguiente tabla el desempeño y las actividades realizadas de cada colaborador, incluyendo a los estudiantes. En caso de cambios en la participación de los investigadores indicarlo, describiendo el grado de afectación al proyecto (máximo media cuartilla).


\begin{table}[h!]
\renewcommand{\arraystretch}{1.3}
\caption{\textsc{Participación e integración del grupo de trabajo}}
\label{tb:grupo}
\centering
\begin{tabular}{lccc}
\bfseries Nombre del  & \bfseries &\bfseries Desempeño en el&\bfseries \\
\bfseries participante & \bfseries Institución &\bfseries Actividades realizadas& \bfseries proceso\\
\hline
Residencias profesionales & 800& 40.03 & 185.07\\
Tesis de licenciatura &2& 0 & 2\\
Trabajo de especialidad& -& - & -\\
Tesis de maestría& -& - & -\\
Tesis de doctorado& -& - & -\\
Publicaciones& - & - & 1\\
Patentes& - & - & 1\\
Prototipos& - & - & 1\\
Asistencia a eventos académicos &- &-&-\\
\hline
\end{tabular}
\end{table}

\subsection{Desviaciones y/o modificaciones}
%1.	Mencione si se hicieron desviaciones y/o modificaciones en el proyecto e indíquelas. Justifique estos cambios. Describa las acciones derivadas de las desviaciones y/o modificaciones. Indique si afectan a los objetivos, metas y metodología planteados en el proyecto original (máximo una cuartilla).

\subsection{Vinculación}
%1.	Si el proyecto tiene vinculación con algún sector, indique las acciones realizadas. 

\subsection{Ejercicio de los montos de apoyo}
%1.	Indique el ejercicio de los montos totales asignados por capítulo.
\begin{table}[h!]
\renewcommand{\arraystretch}{1.3}
\caption{\textsc{Ejercicio del recurso}}
\label{tb:grupo}
\centering
\begin{tabular}{lcc|cc}  
& \bfseries TecNM & &\bfseries Otras fuentes& \\

\bfseries Capítulos & \bfseries Asignado &\bfseries Ejercido& \bfseries Asignado&\bfseries Ejercido\\
\hline
2000 & & & & \\
3000 & & & & \\
Total & & & & \\
\hline
\end{tabular}
\end{table}

\subsection{Problemas o beneficios}
% Manifieste los problemas o beneficios que se han presentado en el desarrollo del proyecto, tanto técnicos como financieros.
\subsection{Anexos}
% Incluya copia de la evidencia de los resultados (publicaciones, constancias, oficio de asignación de director de tesis, actas de examen, registro de trámite de patente, etc.).

%En el caso de contar con información relevante del proyecto, incluirla.

%Disco con el archivo electrónico del Informe Técnico de Avance 
\section{Firma del reporte}
%----------
\begin{center}
\begin{tabular}{ccc}
\bfseries Director del & \bfseries Subdirector  & \bfseries Director del Instituto \\
\bfseries proyecto& \bfseries Académico & \bfseries Tecnológico de \\
\bfseries &\bfseries &\bfseries Morelia\\
\bigskip\\
\bigskip\\
\rule{3.8cm}{0.5pt}&\rule{3.8cm}{0.5pt}&\rule{3.8cm}{0.5pt}\\
\bfseries M.C. Gerardo Marx &\bfseries Dr. Rafael Lara Hernández &\bfseries Ing. Paulino Alberto \\
\bfseries Chávez Campos& &\bfseries Rivas Martínez\\
\end{tabular}
\end{center}
\bibliographystyle{ieeetr}
%\bibliography{Referencias_marx}
%Deshabilita el comentario para que agregues tus referencias


\end{document}